\documentclass[conference]{IEEEtran}
\usepackage[utf8]{inputenc}
\usepackage[spanish,es-tabla]{babel}
\usepackage{amsmath,amssymb,amsfonts}
\usepackage{algorithmic}
\usepackage{graphicx}
\usepackage{textcomp}
\usepackage{xcolor}
\usepackage{cite}
\usepackage{array}
\usepackage{url}
\usepackage[hidelinks]{hyperref}

\title{\textbf{Optimización de la gestión de recursos críticos hospitalarios: integración de pronósticos epidemiológicos y programación lineal ante brotes estacionales}\\[0.5em]
\large Trabajo Final Integrador}

\author{
    \IEEEauthorblockN{Bellor Maria E., Diaz Dylan S., Jimenez Fernando E., Veliz Ignacio M., Zato Sosa Celina A.}
    \IEEEauthorblockA{Universidad Tecnológica Nacional, Facultad Regional Tucumán\\
    Investigación Operativa\\
    \textbf{Profesores:} Dra. Gramajo Guadalupe, Ing. Rojas Cristina, Ing. Fraga Alvaro\\
    Tucumán, Julio de 2026}
}

\begin{document}

\maketitle

\begin{abstract}
    Este trabajo propone un sistema de soporte de decisiones mediante Investigación Operativa para optimizar la logística en hospitales públicos de Tucumán ante brotes estacionales de dengue, bronquiolitis e influenza. La metodología consta de dos fases: primero, un modelo predictivo (Suavización Exponencial Triple Holt-Winters) para proyectar la demanda semanal según tendencia y estacionalidad histórica. Segundo, un modelo de Programación Lineal Entera Mixta Multiobjetivo que optimiza la asignación semanal de siete recursos críticos (camas, UTI, insumos y personal). El algoritmo resuelve la tensión entre minimizar costos/penalizaciones y maximizar la atención priorizando casos de alta letalidad, garantizando eficiencia operativa y rigor ético-médico bajo condiciones de alta presión epidemiológica.
\end{abstract}

\begin{IEEEkeywords}
    Investigación Operativa, Programación Lineal Multiobjetivo, Pronósticos Epidemiológicos, Holt-Winters, Asignación de Recursos Hospitalarios.
\end{IEEEkeywords}

\section{Fundamentación}
Los sistemas de salud pública en Tucumán enfrentan tensiones extremas ante la coexistencia estacional de dengue, bronquiolitis e influenza, provocando la frecuente saturación hospitalaria y el desabastecimiento de insumos críticos \cite{ref1}. Esta volatilidad epidemiológica exige pasar de una respuesta reactiva a una planificación sanitaria proactiva.

Para evitar el colapso, es vital anticipar la demanda de pacientes. El análisis de series temporales mediante la Suavización Exponencial Triple de Holt-Winters (variante Multiplicativa) ha demostrado ser altamente efectivo en la predicción de brotes infecciosos. Su elección se justifica en que la amplitud de los contagios masivos estacionales de dengue, bronquiolitis e influenza crece proporcionalmente al nivel general de la epidemia, permitiendo capturar de manera precisa las tendencias y los marcados patrones estacionales \cite{ref2, ref3}.

Sin embargo, este pronóstico debe acompañarse de una gestión logística eficiente. La aplicación de la Investigación Operativa en sistemas de salud, mediante modelos de programación lineal multiobjetivo, permite optimizar la asignación de recursos escasos, mejorando la utilización de la infraestructura y reduciendo costos operativos al anticiparse a la saturación \cite{ref4, ref5}.

Por lo tanto, este trabajo se fundamenta en la sinergia de ambas herramientas. Los pronósticos de Holt-Winters estimarán la demanda futura, y un modelo de Programación Lineal Entera Mixta optimizará la asignación semanal de recursos hospitalarios finitos, priorizando la atención de los casos de mayor gravedad clínica para garantizar equidad, eficiencia y minimizar la mortalidad.

\section{Objetivos}
Ante la histórica ineficiencia distributiva del sistema sanitario y la consecuente saturación hospitalaria durante picos epidemiológicos, resulta imperativo adoptar herramientas de planificación proactiva que garanticen una atención equitativa y oportuna \cite{ref6}. En respuesta a esta necesidad, se define:

\subsection{Objetivo Principal}
Desarrollar un sistema de soporte de decisiones fundamentado en modelos de Investigación Operativa para optimizar la gestión logística y la asignación de recursos críticos en hospitales públicos de la provincia de Tucumán frente a la demanda estacional aguda por brotes de dengue, bronquiolitis e influenza.

\subsection{Objetivos Específicos}
\begin{itemize}
    \item Proyectar la demanda semanal de atención médica para las patologías en estudio mediante el análisis predictivo de series temporales, implementando el método de Suavización Exponencial Triple (Holt-Winters Multiplicativo) para capturar tendencias y estacionalidades.
    \item Formular un modelo de Programación Lineal Entera Mixta (PLEM) Multiobjetivo de carácter estático que sistematice la asignación de capacidades limitadas, incluyendo camas de internación, unidades de cuidados críticos (UTI), terapias de oxígeno, kits médicos y horas de personal.
    \item Minimizar los costos operativos y las penalizaciones por faltantes de recursos a través de la optimización del algoritmo matemático.
    \item Maximizar la cantidad de pacientes atendidos integrando un factor de ponderación clínica que priorice sistemáticamente los casos de mayor letalidad, garantizando el cumplimiento de los protocolos de resguardo ético-médico.
\end{itemize}

\section{Marco Conceptual}
\begin{itemize}
    \item \textbf{Patología epidémica estacional:} Enfermedad infecciosa aguda con incrementos masivos y cíclicos anuales (dengue, bronquiolitis e influenza).
    \item \textbf{Demanda epidemiológica aguda:} Picos repentinos de pacientes que amenazan con saturar la capacidad instalada y agotar los presupuestos.
    \item \textbf{Recurso crítico asistencial:} Elemento finito e indispensable (camas de internación, UTI, oxigenoterapia, kits médicos y personal de salud).
    \item \textbf{Capacidad neta hospitalaria:} Disponibilidad operativa real, sujeta a restricciones estrictas de presupuesto, espacio y protocolos.
    \item \textbf{Planificación sanitaria proactiva:} Estrategia para anticipar necesidades futuras y evitar el colapso hospitalario frente a crisis epidemiológicas.
    \item \textbf{Ponderación clínica (Triage):} Priorización matemática de atención basada en la severidad y la letalidad histórica de cada paciente.
    \item \textbf{Penalización por déficit de atención:} Costo socioeconómico, legal y ético asumido por el sistema público ante cada paciente no atendido.
    \item \textbf{Logística de Contingencia Sanitaria:} Abastecimiento y distribución eficiente de recursos para sostener la operatividad ininterrumpida.
    \item \textbf{Costo Operativo Hospitalario:} Erogación financiera base y extraordinaria requerida para sostener la infraestructura y el capital humano.
\end{itemize}

\section{Marco Teórico}
\begin{itemize}
    \item \textbf{Pronóstico:} Técnica analítica de la Investigación Operativa que utiliza datos históricos y modelos matemáticos para proyectar escenarios futuros, reduciendo la incertidumbre del sistema \cite{ref7, ref11}.
    \item \textbf{Suavización Exponencial Triple (Holt-Winters):} Método que descompone y proyecta simultáneamente nivel, tendencia y estacionalidad. Es indispensable para estructurar alertas tempranas en enfermedades estacionales \cite{ref9}.
    \item \textbf{Investigación Operativa (IO):} Disciplina cuantitativa que optimiza el rendimiento de un sistema sujeto a restricciones, proveyendo el soporte analítico para readecuar recursos hospitalarios ante brotes \cite{ref10}.
    \item \textbf{Programación Lineal Entera Mixta (PLEM):} Técnica matemática de optimización sujeta a limitaciones lógicas, donde variables clave (como pacientes o camas) toman valores enteros estrictos \cite{ref8}.
    \item \textbf{Modelo Multiobjetivo:} Variante de optimización para resolver problemas con metas en conflicto directo, buscando un equilibrio entre minimizar costos operativos y maximizar la cobertura médica priorizada \cite{ref11}.
\end{itemize}


\section{Marco Contextual}
\subsection{Contexto Demográfico y Población de Riesgo}
El modelo se enmarca en la provincia de Tucumán. Según datos del Censo Nacional 2022 \cite{ref12}, de sus 1.703.186 habitantes, el 42\% depende de manera exclusiva del subsistema público de salud (SIPROSA). La segmentación etaria es fundamental para estimar la demanda potencial asimétrica ante el brote:
\begin{table}[htbp]
    \caption{Distribución demográfica de riesgo y dependencia del sistema público en Tucumán (basado en \cite{ref12})}
    \label{tab:demografia}
    \centering
    \setlength{\tabcolsep}{3pt}
    \begin{tabular}{|>{\raggedright\arraybackslash}m{1.8cm}|>{\raggedright\arraybackslash}m{1.8cm}|>{\raggedright\arraybackslash}m{1.4cm}|>{\raggedright\arraybackslash}m{2.1cm}|}
        \hline
        \multicolumn{1}{|>{\centering\arraybackslash}m{1.8cm}|}{\textbf{Grupo Poblacional}} & \multicolumn{1}{>{\centering\arraybackslash}m{1.8cm}|}{\textbf{Patología de Mayor Riesgo}} & \multicolumn{1}{>{\centering\arraybackslash}m{1.4cm}|}{\textbf{Población Total (Tucumán)}} & \multicolumn{1}{>{\centering\arraybackslash}m{2.1cm}|}{\textbf{Demanda Potencial (Dep. Pública $\approx$42\%)}} \\
        \hline
        Primera Infancia (0 a 4 años)                                                       & Bronquiolitis (VSR)                                                                        & $\approx$ 142.500 hab.                                                                     & $\approx$ 59.850 lactantes y niños                                                                              \\
        \hline
        Población General (15 a 64 años)                                                    & Dengue                                                                                     & $\approx$ 1.085.000 hab.                                                                   & $\approx$ 455.700 jóvenes y adultos                                                                             \\
        \hline
        Adultos Mayores (65 años y más)                                                     & Influenza A / Neumonía                                                                     & $\approx$ 186.200 hab.                                                                     & $\approx$ 78.200 adultos mayores                                                                                \\
        \hline
    \end{tabular}
\end{table}

\subsection{Escenario de Máxima Tensión Epidemiológica (Referencia Histórica)}
Para parametrizar el comportamiento de la demanda ($D_i$) en crisis, se toma el escenario histórico del ciclo 2023-2024 en Tucumán, donde confluyeron el mayor brote histórico de dengue y un aumento temprano de Infecciones Respiratorias Agudas Bajas (IRAB). La saturación máxima se registró entre las Semanas Epidemiológicas (SE) 13 y 17 del 2024:
\begin{table}[htbp]
    \caption{Registro histórico del pico de crisis multi-epidémica en Tucumán (valores de referencia basados en \cite{ref13})}
    \label{tab:pico_critico}
    \centering
    \setlength{\tabcolsep}{3pt}
    \begin{tabular}{|>{\raggedright\arraybackslash}m{1.7cm}|>{\raggedright\arraybackslash}m{1.3cm}|>{\raggedright\arraybackslash}m{2.0cm}|>{\raggedright\arraybackslash}m{2.0cm}|}
        \hline
        \multicolumn{1}{|>{\centering\arraybackslash}m{1.7cm}|}{\textbf{SE (Pico Crítico)}} & \multicolumn{1}{>{\centering\arraybackslash}m{1.3cm}|}{\textbf{Casos Dengue}} & \multicolumn{1}{>{\centering\arraybackslash}m{2.0cm}|}{\textbf{Bronquiolitis ($<$2 años)}} & \multicolumn{1}{>{\centering\arraybackslash}m{2.0cm}|}{\textbf{ETI / Influenza}} \\
        \hline
        SE 13 (Fin Marzo)                                                                   & 7.290                                                                         & 64                                                                                         & 354                                                                              \\
        \hline
        SE 14 (Inic. Abril)                                                                 & 8.333                                                                         & 69                                                                                         & 375                                                                              \\
        \hline
        SE 15 (Med. Abril)                                                                  & 8.160                                                                         & 72                                                                                         & 328                                                                              \\
        \hline
        SE 16 (Fin Abril)                                                                   & 6.878                                                                         & 117                                                                                        & 437                                                                              \\
        \hline
        SE 17 (Inic. Mayo)                                                                  & 6.708                                                                         & 114                                                                                        & 531                                                                              \\
        \hline
    \end{tabular}
\end{table}

Para alimentar y calibrar el modelo predictivo en este contexto de crisis, se definieron criterios metodológicos específicos para la delimitación de las ventanas de datos históricos de cada patología:
\begin{itemize}
    \item \textbf{Dengue (Ciclos 2023--2024):} Se limitó a este periodo debido a que representó el evento de máxima crisis sanitaria en la historia de Tucumán. Incorporar años previos de baja circulación endémica distorsionaría la estacionalidad y restaría capacidad de alerta temprana ante brotes agudos y de comportamiento vectorial altamente dependiente de cepas e indicadores climáticos.
    \item \textbf{Bronquiolitis e Influenza (Ciclos 2023--2025):} Se empleó una ventana de tres años. Estas patologías presentan un comportamiento endoepidémico invernal regular y cíclico, por lo que una muestra más amplia asegura la estabilidad estadística necesaria para modelar las Infecciones Respiratorias Agudas Bajas (IRAB).
\end{itemize}

\subsection{Parámetros Clínicos: Hospitalización, Severidad y Letalidad}
La conversión de casos en requerimientos de insumos emplea tasas de consumo tecnológico ($T_{i,r}$) y ponderadores clínicos de gravedad ($\mu_i$), estructurados bajo directrices de la Organización Mundial de la Salud (OMS) y la Organización Panamericana de la Salud (OPS). El modelo filtra la demanda bruta según tasas de hospitalización esperadas e índices de letalidad:
\begin{table}[htbp]
    \caption{Indicadores clínicos de severidad y requerimientos asistenciales (basado en \cite{ref14, ref15})}
    \label{tab:indicadores_clinicos}
    \centering
    \setlength{\tabcolsep}{3pt}
    \begin{tabular}{|>{\raggedright\arraybackslash}m{1.5cm}|>{\raggedright\arraybackslash}m{1.8cm}|>{\raggedright\arraybackslash}m{1.8cm}|>{\raggedright\arraybackslash}m{2.1cm}|}
        \hline
        \multicolumn{1}{|>{\centering\arraybackslash}m{1.5cm}|}{\textbf{Patología ($i$)}} & \multicolumn{1}{>{\centering\arraybackslash}m{1.8cm}|}{\textbf{Hosp. (Camas)}} & \multicolumn{1}{>{\centering\arraybackslash}m{1.8cm}|}{\textbf{Cuidados Críticos (UTI/ARM)}} & \multicolumn{1}{>{\centering\arraybackslash}m{2.1cm}|}{\textbf{Prioridad / Letalidad ($\mu_i$)}} \\
        \hline
        1. Dengue                                                                         & 5\% - 7\% (Crítica / Alarma)                                                   & $<$ 1\% (Dengue Grave)                                                                       & Bajo / Mod. (Prioridad: Hidratación)                                                             \\
        \hline
        2. Bronquiolitis (VSR)                                                            & 10\% - 15\% (Hipoxemia)                                                        & 3\% - 5\% (UTI)                                                                              & Alto (Riesgo pediátrico rápido)                                                                  \\
        \hline
        3. Influenza A / ETI                                                              & 8\% - 12\% (Riesgo)                                                            & 4\% - 6\% (Insuf. Resp.)                                                                     & Muy Alto (Mortalidad adultos mayores)                                                            \\
        \hline
    \end{tabular}
\end{table}

\subsection{Operativización Financiera y Restricciones de Capacidad}
Los efectores de tercer nivel del SIPROSA operan con capacidad instalada rígida. El modelo calcula la capacidad neta deduciendo la carga base histórica ocupada por patologías no epidémicas (entre 60\% y 70\% de camas). Por su parte, los costos unitarios ($C_r$) y el presupuesto extraordinario semanal ($PresupuestoMax$) se ingresan dinámicamente según valores vigentes de mercado y asignaciones ministeriales, garantizando la aplicabilidad temporal del modelo al calcular la frontera de eficiencia óptima.

\subsection{Origen de Datos y Metodología de Digitalización}
Los datos cuantitativos semanales reales ($y_t$) se obtuvieron de los reportes oficiales de la sala de situación epidemiológica del SIPROSA y de los Boletines Epidemiológicos Nacionales, cubriendo desde marzo de 2023 hasta enero de 2026.

Debido a que múltiples reportes históricos solo presentaban las curvas epidemiológicas en formato gráfico no editable (PDFs), se utilizó la herramienta digitalizadora \textbf{WebPlotDigitizer} para realizar una calibración de ejes cartesianos y extraer con precisión milimétrica las coordenadas numéricas de casos semanales reales.

\section{Modelo simbólico}
El desarrollo del sistema integrado se compone de dos fases matemáticas sucesivas: un modelo predictivo basado en series temporales para estimar la demanda y un modelo de programación lineal entera mixta para optimizar la asignación de recursos.

\subsection{Modelo de Pronósticos (Holt-Winters Multiplicativo)}
La fase predictiva opera sobre un conjunto estructurado de índices y parámetros, donde $t$ denota la semana epidemiológica ($t = 1, 2, \dots, n$), $i$ representa la patología ($i = 1$ para Dengue, $i = 2$ para Bronquiolitis y $i = 3$ para Influenza), $s$ es la longitud del ciclo estacional anual ($s = 52$ semanas) y $m$ es el horizonte temporal de pronóstico.

El modelo de Suavización Exponencial Triple (Holt-Winters Multiplicativo) descompone la serie temporal real ($y_t$) actualizando semanalmente sus componentes a través de constantes de suavizamiento optimizadas en el intervalo $[0,1]$: de nivel ($\alpha$), de tendencia ($\beta$) y de estacionalidad ($\gamma$).

La ecuación de actualización del nivel base $L_t$ se formula como:
\begin{equation}
    L_t = \alpha \frac{y_t}{S_{t-s}} + (1 - \alpha) (L_{t-1} + T_{t-1})
\end{equation}
donde $L_t$ representa el nivel base, $y_t$ el valor real observado y $S_{t-s}$ el factor estacional del año previo. La tendencia secular $T_t$ se actualiza como:
\begin{equation}
    T_t = \beta (L_t - L_{t-1}) + (1 - \beta) T_{t-1}
\end{equation}
El factor de estacionalidad multiplicativo $S_t$ se calcula mediante:
\begin{equation}
    S_t = \gamma \frac{y_t}{L_t} + (1 - \gamma) S_{t-s}
\end{equation}
Finalmente, la proyección de la demanda esperada para un horizonte futuro $m$ se formula como:
\begin{equation}
    D_i = F_{t+m} = (L_t + m T_t) S_{t-s+m}
\end{equation}
Para cuantificar la desviación y precisión de las proyecciones frente a los valores reales, se computan el Error Absoluto ($EA_t$) y la Suma de Errores Cuadráticos ($SEC$):
\begin{equation}
    EA_t = |y_t - F_t|
\end{equation}
\begin{equation}
    SEC = \sum_{t} (y_t - F_t)^2
\end{equation}

\subsection{Modelo de Optimización (Programación Lineal Entera Mixta)}
El modelo PLEM multiobjetivo balancea la sustentabilidad presupuestaria y la efectividad socio-sanitaria en la asignación de recursos críticos.

\subsubsection{Funciones Objetivo}
\begin{itemize}
    \item \textbf{Minimización de Costos Operativos y Penalizaciones ($Z_1$):} Modela el costo total extraordinario semanal asociado al brote. Busca minimizar la suma de la asignación e insumos de recursos adicionales (ponderada por sus costos unitarios $C_r$) y el impacto/penalidad punitiva, socio-ética y legal ($P_i$) incurrida por cada paciente de patología $i$ no atendido ($F_i$).
    \item \textbf{Maximización de la Atención Sanitaria Ponderada ($Z_2$):} Garantiza una priorización ética-médica (triage algorítmico) maximizando la cantidad de pacientes atendidos ($X_i$), ponderados individualmente por el índice clínico de gravedad o letalidad relativa ($\mu_i$) de su patología.
\end{itemize}

\subsubsection{Parámetros del Modelo}
\begin{itemize}
    \item \textbf{Tasa de consumo tecnológico ($T_{i,r}$):} Coeficiente técnico que define los requerimientos específicos del recurso $r$ por paciente de patología $i$.
    \item \textbf{Capacidad nominal ($K_{\text{max},r}$):} Disponibilidad nominal total del recurso $r$ en la infraestructura física y de personal del hospital.
    \item \textbf{Carga base hospitalaria ($B_r$):} Consumo del recurso $r$ ocupado de forma ineludible por pacientes traumatológicos, oncológicos, cardiovasculares y otras patologías habituales no epidémicas.
    \item \textbf{Costo unitario del recurso ($C_r$):} Costo unitario extraordinario por unidad del recurso $r$ activado o consumido.
    \item \textbf{Penalización por desatención ($P_i$):} Valor de penalidad ética, social y legal por cada paciente no atendido de la patología $i$.
    \item \textbf{Ponderador clínico de gravedad ($\mu_i$):} Índice de letalidad o criticidad de la patología $i$ basado en el triage de la OMS/OPS.
    \item \textbf{Presupuesto disponible ($Presupuesto$):} Fondo financiero máximo extraordinario asignado semanalmente para la contingencia sanitaria.
\end{itemize}

\section{Resolución}

\subsection{Fase Predictiva: Resultados y Ajuste del Pronóstico}
Para hallar la combinación óptima de las constantes de suavizamiento ($\alpha, \beta, \gamma$), se empleó el Solver de Microsoft Excel \cite{refExcel}, configurando como celda objetivo la minimización de la Suma de Errores Cuadráticos ($SEC$) sobre los datos históricos.

El modelo de Holt-Winters Multiplicativo logró capturar las dinámicas específicas de cada patología:
\begin{itemize}
    \item \textbf{Dengue:} Mostró un comportamiento altamente explosivo y concentrado en verano y otoño, con factores estacionales muy marcados que multiplican el nivel base, típico de un brote vectorial agudo.
    \item \textbf{Bronquiolitis:} Evidenció picos predecibles durante los meses de invierno con baja dispersión interanual, facilitando la planificación de camas pediátricas y kits de ventilación respiratoria.
    \item \textbf{Influenza:} Reveló una circulación viral sostenida a lo largo del año con picos invernales superpuestos, permitiendo una proyección estable útil para prever horas de personal de salud y consumo de oxígeno.
\end{itemize}

A continuación se presentan los gráficos que comparan las series históricas reales de casos registrados contra las proyecciones obtenidas mediante el modelo predictivo calibrado para cada una de las patologías bajo estudio:

\begin{figure}[htbp]
    \centering
    \includegraphics[width=\linewidth]{docs/Grafico_Dengue.png}
    \caption{Comparación de casos reales vs. pronóstico de Dengue en Tucumán (2024).}
    \label{fig:pronostico_dengue}
\end{figure}

\begin{figure}[htbp]
    \centering
    \includegraphics[width=\linewidth]{docs/Grafico_Bronquiolitis.png}
    \caption{Comparación de casos reales vs. pronóstico de Bronquiolitis en Tucumán (2024-2025).}
    \label{fig:pronostico_bronquiolitis}
\end{figure}

\begin{figure}[htbp]
    \centering
    \includegraphics[width=\linewidth]{docs/Grafico_Influenza.png}
    \caption{Comparación de casos reales vs. pronóstico de Influenza en Tucumán (2024-2025).}
    \label{fig:pronostico_influenza}
\end{figure}

\subsection{Interfaz de Integración: Tabla de Uso}
Dado que el modelo matemático de asignación de recursos en la fase de optimización es de carácter estático y semanal, se diseñó la herramienta denominada ``Tabla de Uso''. Esta interfaz actúa como un filtro dinámico que permite a los planificadores hospitalarios seleccionar una semana epidemiológica específica de interés (por ejemplo, el pico histórico de la SE 17). La Tabla de Uso extrae la demanda proyectada ($D_i$) del modelo de pronósticos y la transfiere de forma automática como un parámetro constante hacia las restricciones del modelo de Programación Lineal Entera Mixta (PLEM).

\subsection{Fase de Optimización: Resolución del PLEM}
% Esta sección se completará al resolver el modelo matemático PLEM.

\section{Interpretación}

\subsection{Interpretación del Modelo de Pronósticos (Holt-Winters)}
La interpretación del comportamiento de las curvas proyectadas indica que el modelo de Holt-Winters Multiplicativo se adapta eficazmente a la volatilidad de los brotes en la provincia de Tucumán. El análisis del dengue revela coeficientes estacionales sumamente altos concentrados en un intervalo de pocas semanas, lo cual exige una planificación reactiva intensiva y focalizada para evitar la saturación inmediata. Por el contrario, la demanda estacional de bronquiolitis e influenza muestra una distribución más suavizada a lo largo de la temporada invernal, lo que permite estructurar un esquema de abastecimiento preventivo y distribución de turnos médicos programada en el tiempo. La Tabla de Uso permite traducir directamente esta dinámica predictiva en los requerimientos específicos de demanda ($D_i$) para cada escenario de estrés sanitario.

\subsection{Interpretación del Modelo de Optimización Lineal (PLEM)}
% Esta sección se completará una vez resuelto el modelo PLEM y analizados los escenarios de asignación de recursos.

\section{Conclusiones}

\subsection{Conclusiones Parciales: Fase Predictiva}
A partir del desarrollo e implementación de la fase de pronósticos, se concluye que:
\begin{itemize}
    \item La Suavización Exponencial Triple Holt-Winters con estacionalidad multiplicativa es una metodología matemática robusta para modelar brotes epidémicos complejos cuyas variaciones no son lineales ni constantes.
    \item El uso de la herramienta WebPlotDigitizer permitió recuperar información cuantitativa de alta fidelidad a partir de curvas publicadas en reportes oficiales no estructurados, eliminando sesgos en el ingreso de datos históricos.
    \item La Tabla de Uso demostró ser una interfaz funcional efectiva para desacoplar y comunicar secuencialmente un modelo dinámico de series temporales con un modelo estático y discreto de optimización semanal.
\end{itemize}

\subsection{Conclusiones Generales del Trabajo}
% Se completará una vez integrados y evaluados los resultados de la optimización lineal.

\begin{thebibliography}{00}
    \bibitem{ref1} Ministerio de Salud Pública de Tucumán. (2024). Sala de situación epidemiológica. Dirección de Epidemiología.
    \bibitem{ref2} Mahmud, N., Ali, N. A., Pazil, N. S. M., \& Jamaluddin, S. H. (2022). Predicting dengue outbreak in Selangor using Holt-Winters models. International Journal of Academic Research in Business and Social Sciences, 12(1).
    \bibitem{ref3} Mahmud, N., \& Pazil, N. S. M. (2021). Prediction of dengue outbreak: A comparison between ARIMA and Holt-Winters methods. Universiti Teknologi MARA.
    \bibitem{ref4} Shah, D., Marathe, S. A., \& Mehrotra, M. (2023). Operations research applications in healthcare systems: Resource allocation optimization models. Journal of Applied Analytics and Forecasting Research.
    \bibitem{ref5} Shams Eddin, M., \& El Hajj, H. (2025). An optimization-based framework to dynamically schedule hospital beds in a pandemic. Preprints.
    \bibitem{ref6} Dirección de Investigación para la Salud. (2017). Agenda Nacional de Investigación en Salud Pública: Diagnóstico de Situación de Áreas de Vacancia en Investigación Operativa en Salud de Argentina. Ministerio de Salud de la Nación.
    \bibitem{ref7} Anderson, D. R., Sweeney, D. J., Williams, T. A., Camm, J. D., \& Martin, K. (2011). Métodos cuantitativos para los negocios (11.ª ed.). CENGAGE Learning.
    \bibitem{ref11} Winston, W. L. (2005). Investigación de Operaciones: Aplicaciones y algoritmos (4.ª ed.). THOMSON.
    \bibitem{ref9} Ministerio de Salud de la Nación Argentina. (2023). Estrategia de Gestión Integrada para la Prevención y Control del Dengue. Dirección Nacional de Epidemiología y Análisis de Situación de Salud.
    \bibitem{ref10} Sistema Provincial de Salud de Tucumán (SIPROSA). (2024). Plan de contingencia y lineamientos operativos para la readecuación de recursos físicos ante brotes epidémicos. Ministerio de Salud Pública, Gobierno de Tucumán.
    \bibitem{ref8} Bonomo, F., Durán, G., \& Marenco, J. (2020). Modelos de simulación y optimización para la gestión de recursos hospitalarios ante variaciones de demanda. Instituto de Cálculo, Universidad de Buenos Aires - CONICET.
    \bibitem{ref12} Instituto Nacional de Estadística y Censos [INDEC]. (2023). Censo Nacional de Población, Hogares y Viviendas 2022: Resultados definitivos. Indicadores demográficos por provincia.
    \bibitem{ref13} Ministerio de Salud de la Nación. (2024). Boletín Epidemiológico Nacional N° 702 - SE 17.
    \bibitem{ref14} Organización Mundial de la Salud [OMS]. (2023). Dengue y dengue grave: Datos y cifras.
    \bibitem{ref15} Organización Panamericana de la Salud [OPS]. (2022). Actualización Epidemiológica: Virus Respiratorio Sincitial (VSR) e Influenza en la Región de las Américas.
    \bibitem{refExcel} Veliz, I. M., Bellor, M. E., Diaz, D. S., Jimenez, F. E., \& Zato, C. A. (2026). Planilla de soporte y calibración del modelo de pronósticos: \emph{Optimización de la gestión de recursos críticos hospitalarios}. Disponible en: \url{https://bit.ly/pronosticos_invop}
\end{thebibliography}

\end{document}